% Raw/canonical equivalence proof. Inclusion fragment.
% Requires amsmath, amssymb, amsthm and theorem/lemma/proposition environments.
% This is a written proof, not a proof-assistant verification.

\section{Zero-canonical scalar syntax and exact equivalence}
\label{zeroI:section}

\paragraph{Reference calculus.}
Let $N$ be the pinned upstream native scalar calculus with raw terms
\[
 t::=x\mid 0\mid St\mid I(n,a;u.h),
\]
formulas $(s=t)$, $Q_yB$, and $D_{x,y}B$, and the eight inference patterns
Hyp, $I_0$, $I_S$, equational transport, QE, DE, DI by cases, and QI by
equational progress.  In particular,
\[
 I(0,a;u.h)=a,\qquad
 I(Sn,a;u.h)=h[I(n,a;u.h)/u].
\]
We use the upstream admissibility of weakening, cut, capture-free numerical
substitution, equality algebra, and transport in compound formulas.  None
of the new normalization claims is needed to establish those upstream
facts.  All terms and bodies below are fixed finite syntax; no numerical
variable ranges over descriptions of functions.

\subsection{A normalizer which only contracts zero counters}

For already normalized $n,a,h$, define the smart constructor
\[
 C(n,a;u.h):=
 \begin{cases}
 a,&n\text{ is literally }0,\\
 I(n,a;u.h),&\text{otherwise}.
 \end{cases}
\]
Define $\nu$ on raw syntax by
\[
 \nu x=x,\quad \nu0=0,\quad \nu(St)=S\nu t,\quad
 \nu I(n,a;u.h)=C(\nu n,\nu a;u.\nu h).
\]
A canonical term is a fixed point of $\nu$.  It contains no iterator with
literal zero counter.  The bound variable belongs only to the body, so
returning the initializer when a node contracts never removes a binder
from the initializer's scope.  Extend $\nu$ to formulas and contexts by
normalizing numerical terms and preserving every formula constructor.

\textbf{Important distinction from the incoming canonical-$R$ report.}
Here $\nu$ can change a nonzero outer constructor into literal zero.  The
test at an iterator must therefore inspect $\nu n$, not the original $n$.
For example,
\[
 I(I(0,0;v.0),a;u.h)
 \quad\text{normalizes to}\quad \nu a.
\]
The false assertion $\nu t\equiv0\Longleftrightarrow t\equiv0$ must not be
imported from the successor-seeded setting.

\begin{lemma}[Elementary properties of normalization]
The map $\nu$ is idempotent, commutes with renaming, introduces no new free
variables, and is primitive recursive on codes of finite syntax trees.
It produces a term no larger than its input tree.
\end{lemma}
\begin{proof}
Structural recursion performs the displayed construction.  A contraction
only deletes an iterator node and its counter and body, retaining its
already normalized initializer; otherwise the recursively normalized
subtrees are retained under the same constructor.  Consequently the
result has no zero-counter node and has no more nodes than the input.
Idempotence, renaming, and the free-variable assertion follow by induction.
Structural recursion on finite coded trees, with these size bounds, is a
primitive-recursive syntactic operation.  This procedure never executes
successor-counter iteration: $I(Sn,a;h)$ is retained, even for a closed
standard numeral $n$.  Thus no evaluator for arbitrary primitive-recursive
programs is being treated as a primitive-recursive operation.
\end{proof}

For a capture-free simultaneous substitution $\sigma$, write
\[
 T\langle\sigma\rangle:=\nu(T[\sigma]).
\]
Under a binder, substitution is lifted in the usual capture-free manner
(or by de Bruijn lifting).  Images may first be normalized without changing
the result.

\begin{lemma}[Substitution coherence]\label{zeroI:coherence}
For every raw term or formula $T$ and every capture-free simultaneous
substitution $\sigma$,
\[
 \nu(T[\sigma])=
 \nu\bigl((\nu T)[\nu\circ\sigma]\bigr).
\]
In particular, normalized substitution is associative and has the usual
identity and binder-lifting laws on canonical syntax.
\end{lemma}
\begin{proof}
Induct on $T$, freshening bound names if names rather than indices are
used.  At an iterator, normalize its substituted counter first.  The
counter induction hypothesis says that both sides obtain exactly the
same normalized counter.  If that counter is zero, both return the
same normalized substituted initializer.  Otherwise both retain an
iterator with the same normalized counter, initializer, and body.

There is one apparently exceptional case: the original normalization
may already have deleted the iterator because $\nu n=0$.  The counter
induction hypothesis then says
$\nu(n[\sigma])=\nu(0[\nu\circ\sigma])=0$, so it is deleted after
substitution as well.  Formula constructors are preserved and their
binder cases use the same argument with lifted substitutions.
Associativity follows by applying the equation twice together with
associativity of ordinary capture-free substitution.
\end{proof}

\subsection{Body extensionality is a proof obligation}

The zero normalizer can simplify an iterator's body.  Numerical
congruence for an ordinary argument of a function symbol does not by
itself justify replacing the entire description of its body.  The
following explicit lemma supplies what is needed.

\begin{lemma}[Raw iterator body extensionality]\label{zeroI:body}
Suppose $u$ is absent from $\Gamma,n,a$ and the raw calculus proves
\[
 \Gamma\vdash_N h(u)=k(u).
\]
Then
\[
 \Gamma\vdash_N I(n,a;u.h)=I(n,a;u.k).
\]
The derivation uses induction for a literal equation only.
\end{lemma}
\begin{proof}
Choose $z$ fresh and put
\[
 H_z=I(z,a;u.h),\qquad K_z=I(z,a;u.k),\qquad E(z)=(H_z=K_z).
\]
The base $E(0)$ follows from two zero computations and equality algebra.
In the step, assume $E(z)$.  Substituting $H_z$ for $u$ in the supplied
proof gives $h(H_z)=k(H_z)$ in the same retained context, since $u$ is
absent from $\Gamma$.  Congruence for $k$ applied to $H_z=K_z$ gives
$k(H_z)=k(K_z)$.  The two successor computations then give the equality
path
\[
 I(Sz,a;h)=h(H_z)=k(H_z)=k(K_z)=I(Sz,a;k).
\]
QI by equational progress supplies $Q_zE(z)$; QE with the base proves
$E(n)$.  No induction on a quantified formula is used.
\end{proof}

\begin{lemma}[Normalization is provably sound]\label{zeroI:soundnorm}
For each raw numerical term $t$, with all its free parameters retained,
\[
 \vdash_N t=\nu t.
\]
For each raw native formula $A$, the formulas $A$ and $\nu A$ are
interderivable in every retained context.  In the ordinary first-order
realization this is the provable biconditional $A\leftrightarrow\nu A$.
\end{lemma}
\begin{proof}
Use external induction on finite term syntax.  Variables and zero use
reflexivity; successor uses numerical congruence.  At $I(n,a;u.h)$,
first use the counter and initializer induction hypotheses to replace
$n,a$ by $\nu n,\nu a$.  The body induction hypothesis, with $u$ still
a free numerical parameter, is $h(u)=\nu h(u)$.
Lemma~\ref{zeroI:body} replaces the body.  We have reached
$I(\nu n,\nu a;u.\nu h)$; if $\nu n=0$, one zero computation reaches
$\nu a$, and otherwise this is already $\nu t$.

For the formula assertion, equality atoms use the proved term equalities
and equality algebra.  It remains to lift pointwise interderivability
through the native binders; this must not be confused with capturing a
bound variable in a single equality-transport template.

Suppose $B(y)$ and $B'(y)$ are pointwise interderivable.  Under $Q_yB$,
assume $B'(0)$, translate to $B(0)$, apply QE at $Sz$, and translate
$B(Sz)$ to $B'(Sz)$.  Derived case-introduction for $Q$ gives $Q_yB'$.
This derived rule follows from DI by cases with an unused existential
coordinate, followed by DE whose local goal is its retained $Q$ clause.
Reverse the pointwise translations for the converse implication.

For $D_{x,y}B$, apply DE at a fresh $b$, retaining $Q_yB(b,y)$.  To
introduce $D_{x,y}B'$ at witness $b$, assume $B'(b,0)$, translate to
$B(b,0)$, use QE at $Sz$, and translate to $B'(b,Sz)$.  DI by cases and
then DE finish the derivation.  Reverse the translations for the other
direction.  Applying these two lifting constructions in an external
induction on formula formation proves the assertion.  No arithmetic
induction on a quantified formula is added.
\end{proof}

\subsection{An unrestricted canonical comparison calculus}

Define $C_*$ on canonical syntax using precisely these eight patterns:
Hyp, reflexivity, successor computation, equational transport, QE, DE,
DI by cases, and QI by equational progress.  Every substitution in these
rules is normalized.  Successor computation is
\[
 C(Sn,a;u.h)=h\langle C(n,a;u.h)/u\rangle.
\]
Equational transport has an arbitrary literal-equation template $E(v)$:
from $s=t$ and $E\langle s/v\rangle$, infer
$E\langle t/v\rangle$.  There is no source restriction and the two
introduction modes are separate.  Thus this intermediate calculus is
stable under canonical substitution and does not yet involve the
literal-zero witness test.

\begin{theorem}[Raw/canonical equivalence]\label{zeroI:equivalence}
For every raw finite context $\Gamma$ and raw native formula $A$,
\[
 \Gamma\vdash_N A
 \quad\Longleftrightarrow\quad
 \nu\Gamma\vdash_{C_*}\nu A.
\]
Both directions are effective transformations of finite derivations.
\end{theorem}
\begin{proof}
For the forward direction, normalize every label of the raw derivation.
A zero computation becomes reflexivity.  A successor computation becomes
the displayed canonical successor computation by
Lemma~\ref{zeroI:coherence}.  Equational transport, QE, DI by cases, and
QI by progress are instances of their canonical versions by the same
coherence lemma.  Hyp and DE retain their shapes.  Freshness is preserved
because normalization introduces no free variables.  The normalized
matrix of an equation is still an equation.

For the reverse direction, view canonical labels as raw syntax.  A
normalized instance $B\langle\sigma\rangle$ is interderivable in $N$
with the corresponding ordinary instance $B[\sigma]$ by
Lemma~\ref{zeroI:soundnorm}.  Consequently each canonical logical or
equational inference can be translated by undoing normalization on its
premises, applying the raw rule, and normalizing its conclusion.
Successor computation translates using the raw successor equation,
$I(n,a;h)=C(n,a;h)$ (the zero equation if $n=0$, reflexivity otherwise),
and normalization soundness on the substituted body.  Reflexivity is
already derived in $N$.  Cut composes these finite derivations.
This yields $\nu\Gamma\vdash_N\nu A$; use normalization soundness on
each assumption and the conclusion to obtain $\Gamma\vdash_N A$.
\end{proof}

For a canonical context and goal, the theorem says exactly
$\Gamma\vdash_{C_*}A$ iff $\Gamma\vdash_NA$.  The round trip from
canonical syntax through raw syntax is literally the identity.  The
other round trip is $\nu$ and is internally the identity by
Lemma~\ref{zeroI:soundnorm}.

\subsection{Interface to the six-rule seed calculus}

Section~\ref{seed:main} supplies the equality and introduction reductions
from $C_*$ to its six-rule calculus $K=\mathcal S$. The auxiliary calculus
$\mathcal A$ there has seven rules; $C_*$ additionally lists reflexivity,
which is already derived from IS and equational transport. Thus
$C_*=\mathcal A$ on every canonical sequent.  The required seed rule is
\[
 \frac{s=t\qquad C(n,s;u.h)=r}{C(n,t;u.h)=r}.
\]
It has no source restriction.  At $n=0$ it becomes root replacement
$s=t,\ s=r\;\Longrightarrow\;t=r$.  With
$c_s:I(S0,s;u.u)=s$, two copies of $c_s$ in this root instance give
$s=s$.  A major $s=t$ and minor $s=s$ then give symmetry $t=s$.
A major $t=s$ and minor $t=r$ give transitivity $s=r$.

At $n=1$, write $A_s=I(S0,s;u.h)$ and $A_t=I(S0,t;u.h)$.
Successor computation gives $A_s=h\langle s/u\rangle$ and
$A_t=h\langle t/u\rangle$.  Seed transport along $s=t$ gives
$A_t=h\langle s/u\rangle$.  Now use this last equation as the major,
the second computation as the minor, and the root instance of Seed.
The result is exactly
\[
 h\langle s/u\rangle=h\langle t/u\rangle.
\]
Thus two successor computations and two Seed inferences give every
term-template congruence.  From a supplied $P$-node proof of $s=t$ the
construction has $P+4$ primitive-inference nodes, counting repeated
subproofs separately.  This is not a bound on certificate bytes or
runtime: the labels still contain the chosen template and endpoints.
Equality algebra and congruence recover full equational transport.
In the other direction the seed rule is directly a special case of
equational transport in $C_*$.

The one-introduction rule is the usual witness fold: use progress exactly
when the matrix is an equation and the canonical witness is literally
zero, and cases otherwise.  The closed canonical term
\[
 d:=I(S0,0;u.u)
\]
is not literal zero and has the one-step computation $d=0$.  Therefore
the omitted equation-case instance with zero witness is recovered by
transporting its hypothesis and conclusion between witnesses $0$ and
$d$, then applying the permitted case instance at $d$.  Only atomic
transport is required, since only an equation matrix loses that case.
The unneeded existential coordinate of a progress instance is eliminated
to recover QI by progress.  These reductions establish $K=C_*$ on all
canonical sequents.

The literal-zero witness test is not preserved by direct substitution
into the same inference labels.  Do not prove substitution admissibility
by silently assuming that it is.  Establish the unrestricted equivalence
above first, then transfer admissibility through $K=C_*$.  Weakening and
assumption cut themselves preserve labels and witness tests and have the
usual direct proofs after freshening eigenparameters.

\paragraph{Template-free checking of seed transport.}
If the major equation is $s=t$ and the minor is $b=r$, its conclusion can
be checked without supplying $n,h$, or any equality template in the
certificate.  If $b$ is literally $s$, output $t=r$.  Otherwise require
$b=I(n,s;u.h)$ and output $I(n,t;u.h)=r$.  The latter counter is already
nonzero by canonical syntax.  These alternatives implement exactly the
smart-constructor rule.  A finite acyclic term cannot both equal $s$ and
contain $s$ as its immediate initializer, so the alternatives do not
conflict.  This is a restriction of the equality inference interface,
not omission of a template validator or a free appeal to arbitrary
congruence.

\subsection{Whole-domain bi-interpretability transfer}

The pinned upstream result supplies ordinary whole-domain interpretations
$F:N\to\mathsf{iPRA}$ and $G:\mathsf{iPRA}\to N$, actual equality,
and both internal composite isomorphisms with graph $x=y$.  Let
$e:C_*\to N$ be the raw reading of canonical syntax and
$q:N\to C_*$ be normalization.  The preceding theorem and its term
and formula lemmas prove
\[
 qe=\mathrm{id}\quad\text{literally on canonical syntax},
 \qquad eq\simeq\mathrm{id}\quad\text{internally in }N.
\]
Define the interpretations for $K=C_*$ by
\[
 F_K:=F\circ e,\qquad G_K:=q\circ G.
\]
Both domains are still the entire numerical domain.  Normal forms are
representatives of \emph{expressions}, not a restriction or quotient of
the domain of numbers.  Equality is actual equality throughout.

For a canonical term or formula, the composite is
$qGFe\simeq qe=\mathrm{id}$ by the upstream round trip and
normalization.  In the other direction it is
$FeqG\simeq FG\simeq\mathrm{id}$.  The analogous assertions for term
equalities and arbitrary formula biconditionals follow from the
corresponding upstream assertions and
Lemma~\ref{zeroI:soundnorm}.  Hence both numerical composite graphs
remain $x=y$.  In particular $K$ and full first-order iPRA are
bi-interpretable and equiconsistent.  The designated contradiction
$0=S0$ is unchanged by normalization, so the exact derivability
equivalence also transfers inconsistency directly.

\paragraph{Ordinary first-order reading.}
As upstream, binder notation for $I$ abbreviates a finite dependency
ordered family of first-order function descriptions.  The canonical
notation imposes the equation $I(0,a;h)=a$ as a definitional spelling
convention and uses normalized substitution.  Expanding this convention
gives its ordinary first-order realization.  Lemmas
\ref{zeroI:coherence} and \ref{zeroI:soundnorm} are precisely what makes
the expansion commute with substitution up to provable equality.
The argument does not identify a normalization algorithm with an
internal arithmetic evaluation function and does not postulate a
function sort or a uniform evaluator.

\paragraph{Scope of the claim.}
This establishes a genuine finite inference-interface reduction and a
whole-domain bi-interpretation.  It does not prove absolute minimality,
does not establish any checker byte improvement without actual complete
implementations and a matched comparison, and is a written proof rather
than a new Lean or Rocq development.
